\section{Methodology}
\label{sec:methodology}

The objective of this research is to establish a rigorous, high-precision conversion framework between 
barometric pressure and geometric altitude in complex urban environments. Unlike pure black-box deep learning approaches, 
our methodology decomposes the altitude measurement into a deterministic physical baseline and a learnable spatial correction component. 
To achieve this, we co-design an edge-sensing hardware terminal and a PINF algorithm.

\begin{figure}[htbp]
    \centering
    % Name your image fig2_hardware.png and put it in the image folder
    \includegraphics[width=\columnwidth]{image/figure2.pdf}
    \caption{Hardware architecture and signal flow of the GeoBox data acquisition terminal. The system integrates GNSS, barometric, and environmental sensors. An onboard MCU performs strict temporal synchronization and simple sensor fusion before transmitting the data payload via a 4G LTE/MQTT interface.}
    \label{fig:hardware}
\end{figure}

\subsection{Edge Sensing Instrumentation}
\label{subsec:hardware}
The foundation of our proposed framework is a custom-designed, low-cost multi-sensor data acquisition terminal, designated as the ``GeoBox''. This edge device is engineered to capture synchronized navigation and atmospheric profiles across distributed urban topologies, providing the sparse but high-fidelity data required to train the continuous neural field. The hardware architecture integrates three core sensing modules:
\begin{enumerate}
    \item \textbf{GNSS Positioning Module}: An integrated GNSS receiver (u-blox M10S) coupled with an external active antenna acquires geometric coordinates (latitude $\phi$, longitude $\lambda$) and the Geometric Altitude ($h_\text{GNSS}$). It supports concurrent reception of multiple constellations (GPS L1C/A, BeiDou B1I, Galileo E1) to ensure robust satellite lock and minimize multipath degradation in urban canyons. Typical horizontal accuracy: 1.5\,m CEP; vertical accuracy: ${\sim}2.5$\,m under open-sky conditions.
    \item \textbf{Barometric Pressure Sensor}: A high-resolution MEMS-based static pressure sensor (Murata ZPA4756) provides raw atmospheric pressure measurements ($P_\text{obs}$) with an operating range of 300 to 1100\,hPa, a resolution of 0.02\,hPa (${\sim}0.17$\,m equivalent at sea level), and long-term stability of $\pm0.5$\,hPa/year. Crucially, the sensor is isolated within a specially designed enclosure featuring vented static ports to minimize dynamic pressure artifacts induced by wind gusts or platform motion.
    \item \textbf{Environmental Sensing Array}: A digital hygrometer and thermometer (Sensirion SHT35) with $\pm0.1\,^{\circ}$C temperature accuracy and $\pm1.5\%$ RH humidity accuracy measures ambient temperature ($T_\text{obs}$) and relative humidity ($RH$). These parameters are required to calculate the specific humidity and virtual temperature, compensating for air density variations in the hydrostatic calculations.
\end{enumerate}

\begin{table}[t]
\caption{GeoBox V1 Key Sensor Specifications}
\label{tab:hardware_specs}
\centering
\small
\begin{tabular}{lll}
\toprule
\textbf{Sensor} & \textbf{Parameter} & \textbf{Specification} \\
\midrule
\multirow{2}{*}{u-blox M10S} & Constellations & GPS + BeiDou + Galileo \\
 & Horizontal accuracy & 1.5\,m CEP \\
 & Vertical accuracy & ${\sim}2.5$\,m \\
\midrule
\multirow{2}{*}{Murata ZPA4756} & Operating range & 300--1100\,hPa \\
 & Resolution & 0.02\,hPa (${\sim}0.17$\,m) \\
 & Long-term stability & $\pm0.5$\,hPa/year \\
\midrule
\multirow{2}{*}{Sensirion SHT35} & Temperature accuracy & $\pm0.1\,^{\circ}$C \\
 & Humidity accuracy & $\pm1.5\%$ RH \\
\bottomrule
\end{tabular}
\end{table}


\subsubsection{Data Synchronization and Transmission}
Metrological integrity depends heavily on the temporal synchronization of the heterogeneous data streams. An onboard Microcontroller Unit (MCU) performs strict hardware-level time synchronization, utilizing the GNSS timing signal Pulse-Per-Second (PPS) to ensure that the geometric altitude is precisely time-stamped alongside the barometric and environmental readings. To guarantee data reliability during continuous field deployments, the system executes a Built-In Test (BIT) sequence upon startup to verify GNSS lock status (minimum satellite count and PDOP threshold), network connectivity, and sensor health. Data collection is delayed until the GNSS module achieves 3D fix with a reported accuracy estimate below a configurable threshold. Powered by a lithium-ion battery for extended autonomy, the MCU aggregates the data into a standardized payload comprising a unique device ID (UID), Unix timestamp, $P_\text{obs}$, $T_\text{obs}$, $RH$, and $(\phi, \lambda)$ coordinates. Data is sampled at a rate of 1\,Hz and aggregated into 1-minute averaged packets to suppress high-frequency measurement noise. These packets are then transmitted to a centralized cloud server via a 4G/LTE cellular module using the lightweight MQTT protocol. This periodic, synchronized data pipeline forms the empirical basis for the physical baseline estimation and neural correction learning phases.

\subsubsection{Training vs.\ Inference Roles}
The three sensing modules serve distinct roles during training and operational inference. During training, the GNSS module provides the \textbf{ground truth altitude} $h_\text{GNSS}$ (the supervised label), while the barometric and environmental sensors provide the input features. During operational inference, a new uncalibrated sensor only needs to report its barometric pressure $P_\text{obs}$, approximate horizontal position, rough altitude estimate, and local temperature/humidity. The trained PINF model then computes the physics-derived $P_{\text{bias}}$, predicts the pressure correction $\delta P$, and returns the corrected geometric height. Critically, the inference sensor does \emph{not} require high-precision GNSS altitude---this is the key enabler for zero-shot generalization.



\subsection{Physical Baseline Estimation and Problem Formulation}
\label{subsec:physical_baseline}
To isolate the complex, non-linear microclimate disturbances inherent to urban VLL airspace, we first establish a deterministic physical baseline. This baseline serves as the mean function for our estimation framework, derived from fundamental atmospheric physics and the edge-sensor measurements acquired by the GeoBox.

The foundational relationship between atmospheric pressure $P$ and geopotential height $H$ is governed by the hydrostatic equation:
\begin{equation}
\frac{dP}{dH} = -\rho g \ ,
\label{eq:hydrostatic}
\end{equation}
where $\rho$ is the air density and $g$ is the gravitational acceleration. Standard aviation models (e.g., ISA) assume a static dry-air atmosphere \cite{icao1993manual}. However, to achieve better accuracy, variations in air density caused by moisture must be accounted for. We utilize the onboard environmental sensors to compute the in-situ Virtual Temperature $T_\text{v}$:
\begin{equation}
T_\text{v} = T_\text{obs} (1 + 0.61 q) \ ,
\label{eq:virtual_temp}
\end{equation}
where $T_\text{obs}$ is the measured absolute temperature and $q$ is the specific humidity derived from the measured relative humidity ($RH$) \cite{simonetti2024geodetic}. 

By substituting the ideal gas law ($\rho = P / (R_\text{dry} T_\text{v})$) into \eqref{eq:hydrostatic} and integrating, we obtain the Hypsometric Equation, which defines our physical baseline estimator $\hat{H}_\text{phy}$:
\begin{equation}
\hat{H}_\text{phy} = H_\text{ref} + \frac{R_\text{dry} T_\text{v}}{g_0} \ln\left(\frac{P_\text{ref}}{P_\text{obs}}\right) \ ,
\label{eq:hypsometric}
\end{equation}
where $R_\text{dry}$ is the specific gas constant for dry air, $g_0$ is the standard gravity, and $P_\text{ref}$ and $H_\text{ref}$ are the reference pressure and height, respectively. This step accounts for the primary altitude variance driven by macro-meteorological conditions.
Crucially, $\hat{H}_\text{phy}$ represents the orthometric height. To harmonize this standard atmospheric altitude with standard UAS GNSS coordinate frames, we apply the Geoid Undulation anomaly $N(\phi, \lambda)$ retrieved from high-resolution gravity models (e.g., EGM96/EGM2008) to obtain the baseline geometric altitude $\hat{h}_\text{phy}$:
\begin{equation}
\hat{h}_\text{phy} = \hat{H}_\text{phy} + N(\phi, \lambda) \ .
\label{eq:geoid}
\end{equation}

% We define the unified urban altitude estimation problem as learning a pressure correction field:
% \begin{equation}
% P_\text{corrected} = P_\text{obs} + \delta P_\theta(\mathbf{x}, \mathbf{p}) \ ,
% \label{eq:problem_formulation}
% \end{equation}
% where $\delta P_\theta$ represents the pressure correction predicted by the PINF parameterized by weights $\theta$. The final height is obtained by substituting $P_\text{corrected}$ into the hypsometric equation (Eq.~\ref{eq:hypsometric}).
While $\hat{H}_\text{phy}$ provides a mathematically rigorous baseline, it fundamentally fails to capture high-frequency local disturbances, such as urban heat islands, building wake turbulence, and sensor-specific aerodynamic biases, which cannot be modeled by simple analytical equations. 
Unlike previous data-driven methods that attempt to directly predict a geometric height residual ($\Delta h$), we define the unified urban altitude estimation problem as learning a \textit{pressure correction field}:
\begin{equation}
P_\text{corrected} = P_\text{obs} + \delta P_\theta(\mathbf{x}, \mathbf{p}) \ ,
\label{eq:problem_formulation}
\end{equation}
where $\delta P_\theta$ represents the pressure correction predicted by the PINF parameterized by weights $\theta$. The final unified height $h_\text{pred}$ is obtained by substituting $P_\text{corrected}$ back into the hypsometric equation \eqref{eq:hypsometric}. This paradigm shift, i.e. correcting the physical state variable ($P$) rather than the derived geometric output ($h$), is the mathematical cornerstone that enables our system to decouple uncalibrated hardware offsets from the spatial environmental model.


\begin{figure}[htbp]
    \centering
    % Ensure your assembled image is named fig3_framework.png and placed in the image folder
    \includegraphics[width=1\linewidth]{image/fig3new.pdf}
    \caption{The overall algorithmic framework. The Physical Baseline Branch derives a deterministic geometric height ($\hat{h}_{phy}$) based on the hypsometric equation. Concurrently, the Neural Correction Branch leverages multi-resolution hash encoding and physics-derived pressure bias features to learn the complex pressure correction field ($\delta P$). The final prediction applies the hypsometric equation to the corrected pressure, supervised by GNSS ground truth.}
    \label{fig:framework}
\end{figure}

\subsection{Multi-Resolution Hash Encoding and Physics-Derived Features}
\label{subsec:hash_and_features}
To effectively learn the highly non-linear pressure correction field, 
the neural network must be capable of resolving high-frequency spatial variations induced by urban microclimates. 
Standard coordinate-based MLPs inherently suffer from ``spectral bias'', favoring low-frequency functions and struggling to fit sharp spatial gradients. To overcome this, we adapt the multi-resolution hash encoding technique \cite{muller2022instant} into our atmospheric measurement framework.

\subsubsection{Multi-Resolution Spatial Encoding}
Unlike traditional positional encoding that relies on fixed sinusoidal functions, hash encoding maps the normalized 3D spatial coordinates $\mathbf{x} = (x, y, z)$ to trainable feature vectors across multiple resolution scales. For a given resolution level $l \in \{1, 2, \dots, L\}$ with a corresponding grid resolution $N_l$, we define a hash table $\mathcal{H}_l$ that maps spatial indices to continuous feature vectors:
\begin{equation}
\mathbf{e}_l(\mathbf{x}) = \mathcal{H}_l\left(\text{hash}\left(\lfloor \mathbf{x} \cdot N_l \rfloor \right)\right) \ ,
\label{eq:hash_lookup}
\end{equation}
where $\mathbf{e}_l \in \mathbb{R}^F$ is the feature vector with dimension $F=4$, and we utilize $L=16$ independent levels to capture spatial interactions ranging from city-wide barometric gradients down to street-level building wakes.

The grid resolutions follow a geometric progression to ensure a balanced receptive field:
\begin{equation}
N_l = N_{\text{min}} \cdot b^l, \quad b = \left(\frac{N_{\text{max}}}{N_{\text{min}}}\right)^{\frac{1}{L-1}} \ ,
\label{eq:resolution}
\end{equation}
with $N_{\text{min}} = 16$ and $N_{\text{max}} = 1024$. The final spatial encoding $\phi(\mathbf{x})$ is the concatenation of features interpolated from all $L$ levels:
\begin{equation}
\phi(\mathbf{x}) = [\mathbf{e}_1, \mathbf{e}_2, \dots, \mathbf{e}_L] \in \mathbb{R}^{L \times F} \ .
\label{eq:concat_encoding}
\end{equation}
This multi-scale approach provides significant metrological advantages: it enables $\mathcal{O}(1)$ constant-time feature lookup for real-time edge processing while dynamically adapting to the non-uniform spatial distribution of the GeoBox sensor network.


\subsubsection{Physics-Derived Pressure Bias Feature}
\label{subsec:bias_feature}

A critical metrological innovation of our approach is the introduction of a physics-derived pressure bias feature that enables zero-shot generalization to previously unseen, uncalibrated sensors. Rather than relying on learned sensor-specific embeddings, which fail to generalize, we compute a continuous physical bias constraint:
\begin{equation}
P_{\text{bias}} = P_{\text{obs}} - P_{\text{expected}} \ ,
\label{eq:pressure_bias}
\end{equation}
where $P_{\text{expected}}$ represents the theoretical static atmospheric pressure at the sensor's true location. The resultant residual, $P_{\text{bias}}$, mathematically isolates the sensor's lumped error, capturing both the static hardware calibration offset and the localized aerodynamic pressure anomaly. This pressure bias is subsequently encoded via a shallow MLP into an 8-dimensional feature vector:
\begin{equation}
\mathbf{f}_{\text{bias}} = \text{MLP}_{\text{bias}}(P_{\text{bias}}) \in \mathbb{R}^8 \ .
\label{eq:bias_encoding}
\end{equation}
This physics-derived formulation marks a fundamental departure from conventional data-driven calibration. By reformulating the neural network's input to rely on an interpretable physical discrepancy rather than arbitrary sensor parameters, the architecture generalizes instantly to new edge devices. The network learns the spatial distribution of the environmental pressure correction independent of the specific hardware collecting the data, directly enabling the exceptional LOSO performance demonstrated in our evaluations.

\subsection{Altitude-Based Curriculum Learning and Network Training}
\label{subsec:curriculum_learning}
Training continuous Neural Fields on spatially heterogeneous IoT data presents severe convergence challenges. 
In an urban VLL environment, measurements acquired near ground level exhibit strong spatial correlation, making them mathematically ``easier'' to fit. 
Conversely, measurements from relatively high-altitude structures are sparse and often suffer from microclimate extrapolation errors, making them ``harder'' samples. To prevent the network from overfitting to densely sampled ground-level data while failing at higher altitudes, we propose an Altitude-Aware Curriculum Learning strategy.

\subsubsection{Three-Stage Curriculum Strategy}
Inspired by cognitive learning principles \cite{bengio2009curriculum}, the training process is structured to expose the network to progressively more challenging geospatial distributions based on altitude:
\begin{itemize}
    \item \textbf{Stage 1 Low Altitude}: Focuses on low-altitude regions to establish fundamental spatial representations. The training subset is defined as $\mathcal{D}_1 = \{ \mathbf{x} \in \mathcal{D} \mid h_{\text{GNSS}} \le 100\text{m} \}$, comprising approximately 30\% of the data, trained for 30 epochs.
    \item \textbf{Stage 2 Moderate Altitude}: Introduces moderate-altitude samples to refine the hash encoding capabilities. The subset is expanded to $\mathcal{D}_2 = \{ \mathbf{x} \in \mathcal{D} \mid h_{\text{GNSS}} \le 200\text{m} \}$, covering $\sim$84\% of the data, trained for an additional 30 epochs.
    \item \textbf{Stage 3 Full Extrapolation}: Incorporates the complete dataset ($\mathcal{D}_3 = \mathcal{D}$), including challenging high-altitude boundary cases ($h_{\text{GNSS}} > 250$m), forcing the network to generalize the physical priors over 80 epochs.
\end{itemize}

\subsubsection{Network Architecture and Loss Formulation}
Before inference, the multi-scale spatial encoding $\phi(\mathbf{x})$ and the physics-derived bias $\mathbf{f}_{\text{bias}}$ must be fused with temporal and environmental contexts. Because atmospheric pressure exhibits strong diurnal cycles, the 1-minute sampling timestamp $t$ is mapped to a continuous 12-dimensional periodic space via Fourier feature encoding, yielding $\mathbf{f}_{\text{time}}$. Additionally, the measured temperature and humidity are concatenated as a 2-dimensional environmental vector $\mathbf{f}_{\text{env}} = [T_{\text{obs}}, RH]$. The comprehensive feature vector is thus formulated as:
\begin{equation}
\mathbf{f}_{\text{input}} = [\phi(\mathbf{x}), \mathbf{f}_{\text{time}}, \mathbf{f}_{\text{env}}, \mathbf{f}_{\text{bias}}] \in \mathbb{R}^{86} \ .
\label{eq:fused_features}
\end{equation}

This composite vector is passed through a 3-layer SIREN (Sinusoidal Representation Network) MLP to predict the continuous pressure correction $\delta P$:
\begin{equation}
\delta P_\theta = \text{SIREN}(\mathbf{f}_{\text{input}}) \in \mathbb{R}^1 \ .
\label{eq:mlp_output}
\end{equation}
SIREN utilizes periodic sinusoidal activation functions, which are uniquely effective for representing physical coordinate-based signals containing high-frequency structural components. To execute Physics-Informed training, the predicted correction $\delta P_\theta$ must be mapped back to the altitude domain. The final predicted geometric altitude $h^{\text{pred}}$ is computed by substituting the corrected pressure ($P_{\text{obs}} + \delta P_\theta$) into the augmented hypsometric equation:
\begin{equation}
h^{\text{pred}}_i = H_{\text{ref}} + \frac{R_{\text{dry}} T_\text{v}}{g_0} \ln\left(\frac{P_{\text{ref}}}{P_{\text{obs},i} + \delta P_\theta}\right) + N(\phi, \lambda) \ .
\label{eq:final_height}
\end{equation}

The model is optimized using the AdamW optimizer \cite{loshchilov2017decoupled} to minimize the MAE between the predicted altitude $h^{\text{pred}}_i$ and the GNSS-derived ground truth $h^{\text{true}}_i$ across a batch size of $B$:
\begin{equation}
\mathcal{L}(\theta) = \frac{1}{B} \sum_{i=1}^B \left| h^{\text{pred}}_i - h^{\text{true}}_i \right| \ .
\label{eq:loss_function}
\end{equation}
Crucially, by formulating the loss on the altitude but predicting the pressure, backpropagation seamlessly computes gradients \textit{through} the physical atmospheric model \eqref{eq:final_height}, enforcing strict physical compliance on the neural weights $\theta$.

To further accelerate convergence, a Cosine Annealing learning rate scheduler with warm restarts \cite{loshchilov2016sgdr} is applied, initialized with a learning rate of $10^{-3}$ and a weight decay of $10^{-5}$. Our ablation experiments confirm that this Physics-Informed Curriculum architecture yields a dramatic 161\% relative improvement, reducing the strict LOSO MAE from 9.27 m (without curriculum) to 3.55 m.



